\scanpage{225}

\mspage{99}

4°) Pour que $(W,C,F)$ soit un triple de Quillen clos et propre à gauche (i.e. $C\subset\mathrm{Cof}_W$) il faut et il suffit qu'on ait la condition (7.4.1.) b) (i.e $W$ contient les morphismes de type injectif), et de plus

\[
C\subset\mathrm{Cof}_W,\quad \unsure{\text{ou}|\text{i.e.}}\inserted{\text{encore :}}\ \text{pour tout carré cocartésien}\tag*{(7.4.2)}
\]

\diagram{\begin{tikzcd}[ampersand replacement=\&,column sep=small,row sep=small]
A \arrow[r,"i"] \arrow[d,dash,"\alpha"'] \& B \arrow[d,dash,"\beta"]\\
A' \arrow[r,"i'"] \& B'
\end{tikzcd}},

avec $i$ un mono, alors $\alpha\in W\Longrightarrow\beta\in W$.

5°) Pour que $(W,C,F)$ soit un triple de Quillen clos et \uline{propre}, il f et s. que les deux conditions de 4°) soient satisfaites, ainsi que l'on ait

\[
F\subset\mathrm{Fib}_W\quad\unsure{\text{ou}|\text{\deleted{ssi}}}\ \text{encore : pour tout carré cartésien}\tag*{(7.4.3)}
\]

\diagram{\begin{tikzcd}[ampersand replacement=\&,column sep=small,row sep=small]
X \arrow[d,dash,"f"'] \& X' \arrow[l,"q"'] \arrow[d,dash,"f'"]\\
Y \& Y' \arrow[l,dash,"p"]
\end{tikzcd}}, si

$f\in F$ (i.e. $f\in(C\cap W)_*$) et $p\in W$, on a $q\in W$.

Dém de 7.4. Les parties 1°, 2°) 3°) ne sont qu'un résumé de ce qui était établi. Pour 4°, il faut voir que les conditions dites impliquent (7.4.1.) a), Or on \deleted{qui est une conséquence formelle du fait qu'on sait déjà que $(C,\mathrm{TF})$ est un couple de Quillen, et \unsure{i.e.}} $C\subset\mathrm{Cof}_W$, donc $C\cap W\subset\mathrm{Cof}_W\cap W\deleted{\unsure{=|{\subset}}\unsure{\mathrm{Cof}}}=\unsure{W}^{\mathrm{univ}}$ \deleted{\unsure{et}}
% archive sheet number: 224
